\section{Discussion}
\label{sec:discussion}

\subsection{Phonetics vs.\ Semantics: Different Distributional Patterns}

The most striking finding is the qualitative difference between how phonetic and semantic features are distributed across layers.
Semantic and syntactic features typically \emph{build} through the network: syntactic information peaks in middle layers, while semantic information strengthens in later layers \citep{radford2019language}.
Phonetic information, by contrast, is strongest at the embedding boundary and gradually degrades.

This makes intuitive sense.
Phonetic form is a property of the token \emph{itself}---it does not depend on context (except for heteronyms).
The embedding table can encode each token's phonetic profile directly, and the first transformer layer refines these representations slightly before the network shifts its representational capacity toward contextual, semantic features.
The progressive compression of phonetic directions in deeper layers (\tabref{tab:geometry}) is the geometric signature of this tradeoff: as the network builds richer semantic representations, it reuses dimensions that previously carried phonetic distinctions.

\subsection{Comparison to Prior Work}

Our results both confirm and extend the findings of \citet{mclaughlin2025rhyme}.
We confirm that phonetic information is linearly recoverable from token embeddings, consistent with their approximately 96\% probe accuracy.
Our lower absolute accuracy (macro F1 = 0.41 vs.\ their 96\%) reflects methodological differences: we use Ridge classifiers with per-phoneme binary predictions and macro-averaged F1 across all 77 phonemes including rare ones, while they use multi-hot logistic probes evaluated differently.
The key comparison is qualitative: we show that phonetic information does not \emph{improve} with depth, which their work did not test.

Our finding that articulatory features do not cluster in the general residual stream contrasts with \citet{mclaughlin2025rhyme}'s organized vowel chart in phoneme mover head result vectors.
This suggests that articulatory organization may be a property of specialized attention heads rather than of the residual stream as a whole---the general representation encodes phonetic information in a less structured way that specialized circuits then organize for the specific task of rhyme production.

\subsection{The Primary-Pronunciation Bias}

The heteronym results reveal something about how language models store and access pronunciation.
The 92\% vs.\ 58\% accuracy gap between primary and secondary pronunciations suggests that embeddings encode a dominant phonetic profile, and contextual processing can sometimes---but not reliably---override it.
Words where the model succeeds (``read,'' ``lead,'' ``tear'') may have more distinctive syntactic cues distinguishing their pronunciations, while words where it fails (``wind,'' ``bow,'' ``live'') have subtler contextual signals.

This connects to our probing finding: if the embedding already commits to a phonetic profile that does not improve with depth, then resolving heteronyms requires the model to \emph{override} embedding-level information using contextual attention---a more demanding computation than simply propagating what the embedding already encodes.

\subsection{Limitations}

\para{Single-token constraint.}
We analyze only words that GPT-2 encodes as a single token, biasing our sample toward common, short words.
Multi-token words constitute the majority of the vocabulary and may exhibit different phonetic encoding patterns, particularly since their pronunciation must be assembled from subword pieces.

\para{Model and probe choices.}
Our experiments use \gptmed, a smaller and older model than the \llama studied by \citet{mclaughlin2025rhyme}.
The phoneme mover head phenomenon may not exist in GPT-2, or may manifest differently.
Our use of Ridge classifiers rather than logistic regression means absolute probe accuracies are not directly comparable across studies.

\para{Isolated token analysis.}
We extract activations for single tokens without surrounding context, which means we measure what is \emph{stored} in each token's representation rather than what is \emph{computed} via attention from neighboring tokens.
This is intentional---we want to isolate the static phonetic content---but it misses the role of attention mechanisms that \citet{mclaughlin2025rhyme}'s phoneme mover heads exploit.

\para{Heteronym sample size.}
Our context-dependent test uses only 12 heteronyms on a single model (\gptfour), limiting statistical power.
A larger-scale heteronym evaluation across multiple models would strengthen these findings.

\para{Pronunciation ambiguity.}
We use the first transcription from \wikipron for each word.
Some words have multiple valid pronunciations, and our probes may partially reflect this ambiguity.
